%!TEX program = xelatex
\documentclass[12pt]{article}

\usepackage[english]{babel}
\usepackage{xeCJK}

% Set page size and margins
\usepackage[a4paper,top=2cm,bottom=2cm,left=3cm,right=3cm,marginparwidth=1.75cm]{geometry}

% Useful packages
\usepackage{amsmath}
\usepackage{graphicx}
\usepackage[colorlinks=true, allcolors=blue]{hyperref}
\usepackage{xurl}

\usepackage{array}
\newcolumntype{P}[1]{>{\centering\arraybackslash}p{#1}}
\newcolumntype{M}[1]{>{\centering\arraybackslash}m{#1}}

\usepackage{adjustbox}
\usepackage{caption}
\usepackage{subcaption}
\usepackage{siunitx}
\usepackage{float}
\usepackage{booktabs}
\usepackage{tabularx}
\usepackage{multirow}
\usepackage{setspace}

% ----- paragraph -----
\usepackage{titlesec}
\setcounter{secnumdepth}{4}
\titleformat{\paragraph}
{\normalfont\normalsize\bfseries}{\theparagraph}{1em}{}
\titlespacing*{\paragraph}
{0pt}{3.25ex plus 1ex minus .2ex}{1.5ex plus .2ex}
% ----- paragraph ----

\setCJKmainfont{NotoSerifTC-Regular}[
  AutoFakeBold = 3
]
\setCJKmonofont{NotoSerifTC-Regular}

\usepackage{indentfirst}
\setlength\parindent{24pt}
\setlength{\parskip}{0.35em}
\setstretch{1.08}

\begin{document}

\begin{titlepage}
  \begin{center}
      \vspace*{1cm}
       \begin{LARGE} 
           \textbf{
           \\數位電子學實務
           }
       \end{LARGE}

       \begin{LARGE} 
            \textbf{
            \\ -- 期中書面提案報告 --
        }
        \end{LARGE} 

      \vspace{1cm}
        \textbf{抬底仔(á)-Tidy UP}
          
      \vspace{1.5cm}

      \textbf{Group 2: B11303052 蘇盈翰, B11702091 張容爾, B10901144 王俊皓}
       
      \vspace{1.5cm}

      \text{\today}

      \vspace{0.8cm}

      \includegraphics[width=0.7\textwidth]{figure/car.png}
     
      \vfill
          
      \vspace{0.8cm}
     
  \end{center}
\end{titlepage}

\section{探討之議題 Why}

根據國家衛生研究院高齡醫學暨健康福祉研究中心在《臺灣高齡健康與長照服務年報》之統計，台灣由高齡者組成之家戶數逐年攀升，其中獨居長者家戶亦呈現持續成長（國家衛生研究院高齡醫學暨健康福祉研究中心，n.d.），加上台灣已於 2025 年正式邁入超高齡社會，65 歲以上人口占比突破 20\%，高齡生活的安全議題已不容忽視。在眾多高齡風險中，「跌倒」對長者的威脅尤為顯著，依衛生福利部統計處公布之 107 年國人死因統計結果，跌倒高居 65 歲以上長者事故傷害死亡原因第二位（每十萬人 25.7 人）（衛生福利部統計處，2021）。此外，衛生福利部新聞稿並援引國民健康署 106 年「國民健康訪問調查」指出，65 歲以上長者每 6 人就有 1 人在過去一年內曾經跌倒（15.5\%）（衛生福利部，2019）。

值得注意的是，跌倒的成因與居家環境高度相關。國民健康署指出，長者跌倒危險因子可分為社會人口學因子、身心功能與疾病用藥因子，以及環境因子三大類，其中居家環境之照明不足、地板髒亂、地毯未固定與走道堆放雜物等，皆為常見潛在危害（衛生福利部國民健康署，n.d.）。由此可知，長者跌倒不只是「多發生在家中」，家中環境本身的失序更是致跌的關鍵推手。日常生活中，散落的拖鞋、未收好的電線等小型障礙物看似不起眼，卻經常是絆倒長者的主因。對此，國民健康署提出「規律運動、維持居家環境、用藥安全」三大防跌原則（衛生福利部，2019），而本組的提案將選擇從「維持居家環境」切入，期望能透過排除家中常見的地面風險，降低長者因跌倒而受傷的機率。

\section{執行方法之探討 How}
為達成「維持居家環境」的目標，讓長者能安全地在家中走動，我們評估了教育宣導、行為監控、居住環境改造等多種可能路徑，最終選擇以「單一自動化拾物裝置」來直接改善居家地面環境，期望能夠以最低的能耗與空間佔用，讓環境能自動維持安全狀態，同時保留長者最大的活動自由。

以下為我們評估過的方案與取捨理由：

\begin{itemize}
  \setlength{\leftmargin}{2em}
  \setlength{\itemsep}{0.5em}
  \item \textbf{教育宣導、改變長者行為：} Pynoos 等人的研究指出，透過環境改造所達到的防跌效果，往往優於僅依賴行為提醒（Pynoos et al., 2010），單純宣導難以長期落實，故效果有限。
  \item \textbf{影像式跌倒偵測系統：} 此類方法雖能即時發現意外，但相關文獻指出，影像式跌倒偵測的主要挑戰之一即為隱私侵犯，攝影機會捕捉個人外貌、行為與居家環境等敏感資訊，容易引發長者疑慮（Wang et al., 2020）。
  \item \textbf{建議長者遷居至更安全的居住環境：} 研究顯示，遷入長期照護機構本身為高度心理壓力的生活事件，可能會引發焦慮、憂鬱與適應困難（Polacsek \& Woolford, 2022）。為避免過度衝擊長者心理健康，此方案將不列入考慮。
  \item \textbf{聘請專業看護：} 此為目前最常見且有效的方案，但受限於照護人力短缺與成本攀升，\textit{TIME} 報導指出，長照費用對家庭造成顯著經濟負擔，故難以作為普遍可行的解法（TIME, 2025）。
  \item \textbf{較極端觀點 - 以全面健康管理延長壽命：} \textit{Business Insider} 在「The Billionaire Trying to Live Forever」專題中提出以高成本生理監測與生活管理來延壽的案例，但此類方案也是成本極高，缺乏普及可能性（Business Insider, n.d.）。
  \item \textbf{將易絆倒物品設計為自動歸位裝置：} 如 Nissan 於 2018 年推出的「Self-Parking Slippers」將自動歸位技術嵌入拖鞋本身（Nissan, 2018）。然而此設計易受多種不可控因素影響，如拖鞋翻面、位置偏移、感測誤差等，使得高成本的技術未必能帶來相應的實際效益。
\end{itemize}

排除上述方案後，我們將焦點放在以「單一工具」解決居家地面的秩序問題，在不增加長者操作負擔的前提下，自動維持地面的安全狀態。

\section{本組的提案 What}
我們想設計一個「自動撿拾散落物小型機器人」 -- \textbf{抬底仔(á)-Tidy UP}，透過環境類比訊號自主辨識地面遺落物，並使用數位訊號的控制將其拾起並歸位，協助長者維持地面環境秩序與安全。我們的設計架構源自 Cook 與 Das 所提出的智慧環境系統架構（Cook \& Das, 2004），透過「感知、推論、執行」的循環，使環境能主動維持安全條件。

\begin{itemize}
  \setlength{\leftmargin}{2em}
  \setlength{\itemsep}{0.5em}
  \item \textbf{感知：} 結合超音波感測器或人體紅外線感測器，用於測距與避障，即時掌握機器人周遭物體的相對位置，透過影像擷取，應用 YOLO 即時物件偵測模型（Redmon et al., 2016）辨識畫面中的物體屬於目標物（如拖鞋）、障礙物或其他背景。
  \item \textbf{決策：} 以 Arduino 或 Raspberry Pi 作為主要控制元件，結合感測器與影像辨識模組的訊號，進行路徑規劃、目標定位與動作判斷。
  \item \textbf{執行：} 包含馬達驅動系統與拾物構造，當機器人抵達目標物旁後，透過鏟子或夾子將拖鞋拾起，並依預設歸位路徑送回指定位置，恢復地面的秩序。
\end{itemize}

透過此系統，我們希望能在不侵犯長者隱私、不改變長者生活習慣的前提下，降低因小型障礙物而生的絆倒風險，為高齡化社會提供更安全、友善的無障礙地面環境。

\section{所需元件與成本評估 Materials \& Costs}

\subsection{所需元件 Materials}
\begin{table}[H]
    \centering
    \caption{零件清單表}
    \label{tab:materials}
    \scriptsize
    \setlength{\tabcolsep}{3pt}
    \renewcommand{\arraystretch}{0.95}
    \begin{tabularx}{\linewidth}{l X X X}
      \toprule
      車體功能 & 項目 & 用途 & 選擇 \\
      \midrule
      \multirow{3}{*}{感測避障}
        & HC-SR04 超聲波感測器           & 使用超聲波感測前方障礙物 & \multirow{2}{=}{HC-SR04 與 HC-SR501 擇一使用} \\
        & HC-SR501 人體紅外線感測器      & 偵測會發出紅外線的物體的移動 & \\
        & Servo 馬達                   & 使感測器左右旋轉增加範圍 & -- \\
      \midrule
      \multirow{3}{*}{通訊來源}
        & Arduino UNO 主控板              & 為整台拾物車運行的控制來源 & \multirow{3}{=}{Arduino UNO 主控板＋HC-06 藍牙模組 與 Arduino UNO R4 WiFi 擇一使用} \\
        & HC-06 藍牙模組                  & 與外部無線通訊 & \\
        & Arduino UNO R4 WiFi            & 結合上述兩者的功能 & \\
      \midrule
      \multirow{5}{*}{動力\&車身}
        & L298N 馬達驅動板                & 控制馬達的方向與速度 & -- \\
        & DC 馬達                         & 驅動車輪的馬達 & \multirow{4}{=}{DC 馬達＋車輪＋3D 列印底盤 與 現成底盤 擇一使用}   \\
        & 車輪                            & 使車身移動 & \\
        & 3D 列印底盤                      & 客製化車身 & \\
        & 現成底盤\&車輪\&馬達              & 使用現成底版進行客製化 & \\
      \midrule
        能量來源 & 18650 鋰電池             & 一顆供電 3.7 V & -- \\
      \bottomrule
    \end{tabularx}

  \end{table}

\subsection{成本評估 Costs}
\begin{table}[H]
  \centering
  \caption{成本評估表}
  \label{tab:costs}
  \scriptsize
  \setlength{\tabcolsep}{3pt}
  \renewcommand{\arraystretch}{0.95}
  \begin{tabularx}{\linewidth}{l X r c r r}
    \toprule
    車體功能 & 項目 & 單價 & 數量 & 目前選擇 & 最便宜選擇\\
    \midrule
    \multirow{3}{*}{感測避障}
      & HC-SR04 超聲波感測器            & 55  & 1 & -- & 55\\
      & HC-SR501 人體紅外線感測器        & 35  & 4 & 140 & -- \\
      & Servo 馬達                     & 50  & 1 & 50 & 50 \\
    \midrule
    \multirow{3}{*}{通訊來源}
      & Arduino UNO 主控板             & 370  & 1 & -- & 370\\
      & HC-06 藍牙模組                 & 155  & 1 & -- & 155\\
      & Arduino UNO R4 WiFi           & 1{,}200 & 1 & 1{,}200 & --\\
    \midrule
    \multirow{5}{*}{動力 \& 車身}
      & L298N 馬達驅動板                & 40 & 1 & 40 & 40\\
      & DC 馬達                        & 65 & 4 & -- & --\\
      & 車輪                           & 20 & 4 & -- & --\\
      & 3D 列印底盤                     & 約 100 & 1 & -- & --\\
      & 現成底盤\&車輪\&馬達             & 250 & 1 & 250 & 250\\
    \midrule
      能量來源 & 18650 鋰電池            & --  & 2 & 247 & 247\\
    \midrule
    \multicolumn{4}{r}{\textbf{預估總計}} & \textbf{1{,}680} & \textbf{920}\\
    \bottomrule
  \end{tabularx}

\end{table}

\section{預期成果與分工規劃 (When \& Who)}

為確保\textbf{抬底仔(á)-Tidy UP}能如期完成，我們制定了詳細的開發時程，分工部分目前是由張容爾負責領件採買與成本評估、王俊皓負責技術實現與蘇盈翰負責技術統籌，但往後執行將是三個人一起執行當周任務，且再視當周情況做更細節的調整。
下表\ref{tab:schedule_v2}結合了週進度目標與具體執行內容：

\begin{table}[htbp]
    \centering
    \caption{每週預期成果規劃表}
    \label{tab:schedule_v2}
    \renewcommand{\arraystretch}{1.4} % 調整行高，讓表格更易閱讀
    \begin{tabularx}{\textwidth}{c |l| X c} % X 欄位會自動調整寬度並換行
        \toprule
        \textbf{預定日期} & \textbf{階段性目標} & \textbf{預期成果與具體執行項目}  \\
        \midrule
        04/24 & 基礎載具建置 & 完成基礎避障車之硬體組裝，確認底盤移動與馬達控制功能正常。  \\
        05/01 & 基礎感測整合 & 實裝遠紅外線感測模組，完成近距離障礙物偵測與基礎避障邏輯測試。 \\
        05/08 & 進階環境探測 & 導入光學雷達 (LiDAR) 系統進行環境掃描，提升障礙物偵測的精準度。  \\
        05/15 & 回收機構實作 & 完成鏟子初步外型設計與組裝，使機器人具備接觸目標物的能力。 \\
        05/22 & 機構作動優化 & 改良鏟取機構作動流暢度與穩定性，確保物品能確實被回收至基地。  \\
        05/29 & 原型機展示 & 軟硬體系統整合完畢，進行初步可運行成品之進度報告與動態展演。  \\
        06/05 & 系統除錯精進 & 針對測試反饋進行全面除錯、軟體參數調校與硬體細部修整。\\
        06/12 & 期末完整驗收 & 繳交期末報告，並展示具備自主探測與物品回收能力。\\
        \bottomrule
    \end{tabularx}
\end{table}


\section{初步成果樣貌}

\section*{參考文獻}
國家衛生研究院高齡醫學暨健康福祉研究中心。（n.d.）。僅高齡人口居住宅數變化。臺灣高齡健康與長照服務年報。\url{https://ageing.nhri.edu.tw/annual_report/tCSMjW1F27USj12UCMW99}

衛生福利部。（2019 年 9 月 24 日）。每 6 人就有 1 位老人曾跌倒 國健署傳授防跌妙招。\url{https://www.mohw.gov.tw/cp-4253-49428-1.html}

衛生福利部統計處。（2021 年 7 月 2 日）。107 年國人死因統計結果。\url{https://dep.mohw.gov.tw/dos/lp-5069-113-xCat-y107.html}

衛生福利部國民健康署。（n.d.）。老人跌倒常見的危險因子有那些？\url{https://www.hpa.gov.tw/Pages/Detail.aspx?nodeid=807&pid=4327}

Business Insider。（n.d.）。\textit{The Billionaire Trying to Live Forever (Life Extended)}。\url{https://www.youtube.com/watch?v=8W4bLxQAQgQ}

Cook, D. J., \& Das, S. K. (Eds.). (2004). \textit{Smart environments: Technologies, protocols, and applications}. Wiley. \url{https://onlinelibrary.wiley.com/doi/book/10.1002/047168659X}

Nissan. (2018, January 25). \textit{Hotel guests get a kick out of Nissan's self-parking slippers}. \url{https://global.nissannews.com/releases/180125-01-e?lang=en}

Polacsek, M., \& Woolford, M. (2022). Strategies to support older adults' mental health during the transition into residential aged care: A qualitative study of multiple stakeholder perspectives. \textit{BMC Geriatrics, 22}, 186. \url{https://bmcgeriatr.biomedcentral.com/article/10.1186/s12877-022-02859-1}

Pynoos, J., Steinman, B. A., \& Nguyen, A. Q. D. (2010). Environmental assessment and modification as fall-prevention strategies for older adults. \textit{Clinics in Geriatric Medicine, 26}(4), 633--644. \url{https://pubmed.ncbi.nlm.nih.gov/20934614/}

TIME. (2025, July 23). \textit{Why so many seniors can't afford long-term care}. \url{https://time.com/7304510/seniors-long-term-care-costs/}

Wang, X., Ellul, J., \& Azzopardi, G. (2020). Elderly fall detection systems: A literature survey. \textit{Frontiers in Robotics and AI, 7}, 71. \url{https://pmc.ncbi.nlm.nih.gov/articles/PMC7805655/}

Bardaro, G., Antonini, A., \& Motta, E. (2022). Robots for elderly care in the home: A landscape analysis and co-design toolkit. \textit{International Journal of Social Robotics, 14}(3), 657--681. \url{https://link.springer.com/article/10.1007/s12369-021-00816-3}

Kumar, A., et al. (2021). \textit{Autonomous bot with ML-based reactive navigation for indoor environment}. arXiv. \url{https://arxiv.org/pdf/2111.12542}

Redmon, J., Divvala, S., Girshick, R., \& Farhadi, A. (2016). You only look once: Unified, real-time object detection. In \textit{Proceedings of the IEEE conference on computer vision and pattern recognition} (pp. 779--788). \url{https://arxiv.org/abs/1506.02640}

\end{document}